\section{在标记点布设连线}\label{subsec:planting}

输入网络中的一个站点可能与许多其他站点有连线。为了将该站点作为单个插入来移动，我们把这些连线在该站点的各端点编码到一条腿中。此时，球面形式必须能检验每条连线两端的比特是否一致。我们分别在每个网络的每个被占用单元中实现这一检验。

每个单元有两个球面形式。形式 $A$ 承载站点的腿以及通向相邻单元的接缝。$A$ 上的额外标记点提供一个检验连线的联合态。私有形式 $B$ 通过每个额外标记点上单独的一张插入表实现该联合态。因此，布设完毕后，每个插入在每个球面上恰有一条腿，此外还可带有外露比特或控制比特。在原站点和接缝处，$B$ 腿均处于真空态。图~\ref{fig:stack} 展示了这两类插入。

\begin{figure}[htb]\centering
\begin{tikzpicture}[x=1cm,y=1cm,font=\small,
  strip/.style={draw=ink,rounded corners=3pt,line width=.65pt},
  ins/.style={draw=ink,rounded corners=2pt,fill=white,inner xsep=7pt,inner ysep=4pt},
  leg/.style={draw=ink,line width=.8pt}]
\draw[strip,fill=ink!10] (-4.35,.98) rectangle (4.35,1.55);
\draw[strip,fill=ink!5] (-4.35,-1.55) rectangle (4.35,-.98);
\node[anchor=west] at (-4.22,1.26) {缝合球面 \(A\)};
\node[anchor=east] at (4.22,1.26) {\(R\) 个物种};
\node[anchor=west] at (-4.22,-1.26) {私有球面 \(B\)};
\node[anchor=east] at (4.22,-1.26) {\(R/b\) 个物种};
\draw[leg] (-5.35,1.26)--(-4.35,1.26);
\draw[leg] (4.35,1.26)--(5.35,1.26);
\node[above,font=\footnotesize] at (-4.85,1.26) {接缝};
\node[above,font=\footnotesize] at (4.85,1.26) {接缝};
\node[ins] (sta) at (-2.15,0) {原站点};
\node[ins] (aux) at (2.0,0) {辅助插入};
\draw[leg] (sta.north)--(-2.15,.98);
\draw[leg,densely dotted] (sta.south)--(-2.15,-.98);
\node[anchor=west,font=\footnotesize] at (-2.1,-.6) {真空};
\draw[leg] (aux.north)--(2.0,.98);
\draw[leg] (aux.south)--(2.0,-.98);
\node[anchor=west] at (2.05,.62) {\(z^P\)};
\node[anchor=west] at (2.05,-.62) {\(z\)};
\foreach \x/\y in {-2.15/.98,2.0/.98,2.0/-.98}
    \fill[ink] (\x,\y) circle (1.5pt);
\draw[draw=ink,fill=white] (-2.15,-.98) circle (1.5pt);
\end{tikzpicture} \caption{站点表把与该站点相接的内部连线的比特编码在 \(A\) 上。辅助表将 \(A\) 与 \(B\) 耦合，共同施加各连线两端相等的约束。只有 \(A\) 与相邻单元缝合。}\label{fig:stack}
\end{figure}

$B$ 上位于 $z$ 的标记点对应 $A$ 上位于 $z^P$ 的标记点，其中 $P$ 是域的特征的一个幂。将 $z$ 移动 $t$，就使 $z^P$ 移动 $t^P$。因此，以 $A$ 上的洛朗阶为单位时，这一对腿的能量为 $e_A+P^{-1}e_B$。给定同一能量预算，$B$ 上允许的能量是前者的 $P$ 倍。我们在 $B$ 上使用较少的物种，以使对应的腿保持较小的宽度。这将给出站点穿越接缝时留在原处的腿的宽度上界。

按引理~\ref{lem:local} 固定 $g,\ell$。令 $F$ 为不小于 $D^{1000(\ell+1)}$ 的最小特征幂，且令 $a=\lceil F^3\rceil$。给定组数 $M$ 及充分大的特征幂 $b>F$，取 $R$ 为不超过 $T/a$ 的最大 $bM$ 倍数。令
\begin{equation}\label{eq:plant-scales}
 R'=R/b,\qquad P=b/F,\qquad \kappa=P^{-1},\qquad H=Ra^2 .
\end{equation} 其中 $R$ 和 $R'$ 分别是 $A$ 和 $B$ 上的物种数。当 $T$ 充分大时，$T/(2a)\le R\le T/a$。若 $e_A+\kappa e_B\le h$，则由式~\eqref{eq:local-width}，两侧宽度的上界分别为 $\sqrt{Rh}$ 和 $\sqrt{R'Ph}=\sqrt{Rh/F}$ 的常数倍。以比特为单位，忽略取整误差，便有
\begin{equation}\label{eq:plant-widths}
 w_A=O\bigl(T\sqrt{h/H}\bigr),\qquad
 w_B=O\bigl(T\sqrt{h/(FH)}\bigr).
\end{equation}

\subsection{单个物种的连线形式}

将每对不同站点之间的连线合为一束有序线束，把端点处的置换并入站点表。先在这些表内收缩自环。把每条内部线束拆分成每包 $R$ 比特的数据包，末包不足处补零。物种 $s$ 将承载每个包中的第 $s$ 个比特。每个站点与 $O_g(D)$ 束线相接，每个单元连同与之关联的边总共涉及 $O_g(D^2)$ 束线。因此，在元数为 $t_i$ 的站点处，单个物种至多承载 $t_i/R+O_g(D)$ 比特。

选取整数能量 $h_i$ 处不同的楔积基矢来编码站点 $i$ 处的比特串。根据引理~\ref{lem:dimension}，沿整数能级趋于无穷时，单物种空间维数的对数渐近等于 $(\pi/\sqrt3)\sqrt{h_i}$。由于 $R\ge T/(2a)$ 且 $a\ge D$，可以取 \[
                 h_i=O_g\bigl(a^2(1+t_i/T)^2\bigr).
\] 在接缝 $j$ 处，跨过交界面的各包对一个物种共贡献 $O_g(DT/R+D^2)=O_g(Da)$ 比特，故整数能量 $q_j=O_g(D^2a^2)$ 已经足够。两个单元在此处使用相同的比特次序与基矢编码。固定 $D$ 时，每个物种的包数、能量与编码均取自有限集。

在这些站点与接缝的能级空间上，用系数定义多线性形式 $L$。当所有编码均有效且每一比特的两端均一致时，赋予相应基元组系数一。对于通向另一单元的连线，将它的站点端与接缝端相比较。其余系数均为零。我们将在固定的\emph{总}能量处缝合多物种接缝。为恢复每个物种各自正确的能量，令 $Q=\sum_j q_j$，并将 $L$ 在其余满足总和至多为 $Q$ 的非负半整数接缝能量元组上延拓为零。各站点能量保持固定。

\begin{lemma}[权重有界的辅助标记点]\label{lem:plant-probes} 在单物种球面上多项式个额外标记点插入一个联合态，就能实现延拓后的形式 $L$。在每个额外标记点，此态只使用真空与 $\psi_0$。其支撑集内每个码字所占用的标记点数均为偶数且至多为 $O_g(Da)$。\end{lemma}
\begin{proof} 长度为 $d$ 的楔积，其能量至少为 $\tfrac12+\tfrac32+\cdots+(d-\tfrac12)=d^2/2$。各站点处的总长度为 $O_g(a\sum_i(1+t_i/T))=O_g(Da)$。对于至多 $g$ 条接缝，由柯西–施瓦茨不等式及其总能量界 $Q=O_g(D^2a^2)$，也得出相同的长度上界。列出所得输入空间中出现的原标记点与模态的一切数对 $(v,m)$。此列表的大小至多为 $D,a$ 的多项式。

在新点 $z$ 处插入 $\psi_0$ 时，它与上述模态的配对由坐标为 $(z-z_v)^{-m-1}$ 的余向量给出。在任意特征下，这些不同的主部给出线性无关的有理函数。选取互异的正则点，使其余向量张成模态空间的对偶空间。事实上，若某个非零线性组合在所有已选点上均为零，它必在代数闭包的某个未使用的点取非零值。所需点数至多等于列表中的模态数，且这些点都在某一有限扩张内。如有必要可添入真空，使辅助标记点数至少为三。

对由 $d$ 个新标记点构成的一个子集，在这些点上插入 $\psi_0$，在其余点上插入真空。先考虑所得普法夫式中每个辅助生成元均与一个原生成元配对的部分。该部分本身就是至多 $d$ 个标记点上的插入检验的线性组合。事实上，把所有辅助生成元排在末尾，再按每个匹配所含辅助生成元内部配对的确切集合进行分组。每个非空配对集合的贡献等于其收缩所得的标量，乘以一个仅依赖于配对位置的符号，再乘以剩余标记点上没有辅助内部配对的部分。归纳地减去这些贡献，就分离出所需部分。

现取输入空间直和的偶部中的一个非零向量，令 $d$ 为其非零分量的最大楔积总长度。用 $d$ 个标记点且不含辅助内部配对的检验在所有更短的楔积上均为零。在长度为 $d$ 的分量上，它所计算的是选定 $d$ 个余向量的外积在该分量上的值，仅相差一个共同的符号。各个原基元组在以数对 $(v,m)$ 的列表为生成元的外代数中，是互不相同的基楔积。因此，总有一种 $d$ 个余向量的选法能检出该分量。由此得到的偶检验张成偶部的整个对偶空间。实际使用的每个插入都具有不超过 $d=O_g(Da)$ 的偶权重，并在奇向量上取零。规定的整数能量给出偶长度，因此 $L$ 的零延拓在奇部上也取零。插入检验的一个线性组合遂给出所需联合态。\end{proof}

在 $R$ 个物种中各用一份该态。在站点 $i$ 处，插入表将每个物种限制在能级 $h_i$，对数据包解码，检查填充比特为零，并给出带有外露比特的原表项。在接缝 $j$ 处，我们对总能量 $Rq_j$ 处相同的基矢实行缝合。为看出这能恢复原来的连线，将普法夫式按物种分组。先将原标记点和辅助标记点排为前后相接的两个列表，再将每个列表的元素按物种重排。辅助列表重排所引入的符号可并入联合态。由于每种物种的辅助权重均为偶数，将两个列表交错排列不会产生额外符号。原列表重排的符号只依赖输入楔积。在收缩辅助态后所得的任何非零乘积中，由零延拓的定义，每一物种的接缝能量元组或为指定元组，或其和大于 $Q$。所有物种的这些和加起来等于 $RQ$，所以每一物种都必须取指定元组。此时全部原楔积长度均为偶数，剩余的重排符号为一。该乘积恰好检验了每条连线两端的比特是否相等。我们还在每个单元内放置一个真空汇集点，供移动步骤使用。

\subsection{实现辅助联合态}\label{subsec:plant-private}

引理~\ref{lem:plant-probes} 指定了所有辅助标记点共享的一个张量。为了移动，我们需要在每个标记点上分别放置一张插入表。私有球面通过局部映射给出这一共享张量。一个 $B$ 物种将对应 $A$ 上含 $b$ 个物种的一块。

将缝合侧的 $R$ 个物种依次划分成 $R'$ 个连续块。若某块在辅助标记点 $i$ 处有 $d_i$ 个 $\psi_0$ 模态被占用，则称该块的\emph{类型}为 $\tau=(d_i)_{i=1}^m$。将块张量投影到这些局部子空间，得到张量 $J_\tau$，其中计入了块内重排的符号。它的第 $i$ 条腿的维数为 $\binom b{d_i}$。固定由各块类型构成的有序字后，联合张量等于各 $J_\tau$ 的乘积，仅相差一个已知符号。交换来自不同块的两个位点列表，对符号的指数仅贡献两列表长度的乘积模二。只保留满足 $J_\tau\ne0$ 的类型。

用 $H_2$ 表示以比特计的二元熵，其中 $H_2(0)=H_2(1)=0$。取间距为 $1/m$ 的固定正网格，将每个 $H_2(d_i/b)$ 向上舍入到严格大于它的下一个网格点 $u_i$。于是
\begin{equation}\label{eq:plant-ent}
 \log_2\binom b{d_i}\le b u_i,\qquad
 \max_i u_i\le 2,\qquad \sum_i u_i=O_{g,\ell}(Da\log D).
\end{equation} 事实上，由引理~\ref{lem:plant-probes} 有 $s=\sum_i d_i/b=O_g(Da)$。由凹性可知，各熵之和不超过 $s\log_2(em/s)$，并约定在 $s=0$ 时此式取零。固定 $\ell$ 及域的特征时，标记点个数 $m$ 至多为 $D$ 的多项式，且舍入使总和至多增加一。

在对应的 $B$ 标记点上，取所有满足第 $i$ 个位点处能量至多为 $C(bu_i)^2$ 的非零单物种球面形式的直和 $S_\tau$，其中 $C$ 是稍后选取的绝对常数。当 $b$ 充分大时，存在常数 $c_\tau$，使得对每个 $d\ge1$，一组局部线性映射的乘积给出
\begin{equation}\label{eq:plant-comparison}
 \bigoplus_{s=1}^{c_\tau} S_\tau^{\otimes d}
                    \longrightarrow J_\tau^{\otimes d}.
\end{equation} 因为 $c_\tau$ 与 $d$ 无关，同一选择器指标就能供所有出现的该类型共用。

为进行比较，取自然对数，并在代数闭包上工作，那里指数仍严格小于 $9/4$。对位于 $x_1,\ldots,x_m$ 的归一化谱价格，令 $M_u=\inf_o\sum_i u_i\,d(o,x_i)$。在 $o$ 处制备 $J_\tau$，再输运其各条腿，由式~\eqref{eq:plant-ent} 可知，其价格至多为 $b(\log 2)\sum_i u_i d(o,x_i)$。取下确界得到 $b(\log 2)M_u$。对大小为二的对角张量应用相同的输运上界，该张量的价格为 $\log 2$，故得 $\inf_o\sum_i d(o,x_i)\ge1$。因此 $M_u\ge\min_i u_i>0$。取尺度 $N=b$、宽度上限 $c_0\sqrt C\,u_i$，应用推论~\ref{cor:universal}。其中的比较张量恰好是 $S_\tau$，其价格为 $b(c_0\sqrt C\,M_u+o(1))$，且此估计关于归一化价格一致。取 $c_0\sqrt C>2\log 2$。固定 $D$ 时，网格向量和 $A$ 上的位置选择均只有有限种，且均定义在某个固定的有限扩张上。取 $P$ 次根在该域上给出置换，因此当 $b$ 改变时，$B$ 上可能的位置也都落在有限集中。于是，当特征幂 $b$ 充分大时，对每一种价格，$S_\tau$ 都严格支配 $J_\tau$。由第~\ref{sec:prices-background} 节关于域扩张和严格比较的结果，在固定的定义域上得到式~\eqref{eq:plant-comparison}。这些映射可以延拓至输入表所在的域。

固定类型字。对于每种重数为 $d>0$ 的类型，选取一个副本指标 $s$，并为 $d$ 个因子分别选取一个 $B$ 能量元组，依此展开式~\eqref{eq:plant-comparison} 中的限制。按原来的块顺序将各因子分配给私有物种。它们的乘积就是多物种球面形式，仅相差一个置换齐次列表产生的符号。此时，各位点映射给出一张局部表，将该标记点处的 $A$ 腿和 $B$ 腿耦合起来。它可以联合地作用于\emph{该标记点处}同一类型的所有块。对所有选择和类型字求和，便得到所需的辅助张量。原站点和接缝处的 $B$ 腿仍保持真空态。

\paragraph{控制与存储。}一个\emph{私有选择器}记录有序类型字、每一块的 $B$ 能量以及每种出现的类型所对应的一个指标 $s$。用单张表检查这些数据的合法性，包括未用槽位是否补零，并给出置换符号。由于 $d_i\le b$ 且能量的两倍为 $O(b^2)$，每个块记录类型和两倍能量共需 $O(m\log(2b))$ 比特。全部副本指标总计耗费 $O_{D,b}(1)$ 比特。共有 $R/b$ 个块，对任意固定 $\epsilon>0$，先取充分大的 $b$，再取充分大的 $T$，就可使选择器的长度小于 $\epsilon T$。

在 $A$ 上，每个被占用的辅助模态耗费能量 $1/2$，故辅助总能量为 $O_g(RDa)$。在一个 $B$ 物种上，由式~\eqref{eq:plant-ent}，各能量上限的和为 $O(b^2\sum_i u_i^2)=O_{g,\ell}(b^2 Da\log D)$。乘以 $R'$ 并换算为 $A$ 的单位，得到每个单元中的能量界 \[
 \sum_i(e_{A,i}+\kappa e_{B,i})
    =O_{g,\ell}(RDa+F RDa\log D)\le D H
\] 这里取 $D$ 充分大，并用到了 $H=Ra^2$ 与 $a\ge F^3$。就存储而言，使用所选块的实际空间。在标记点 $i$ 处，由式~\eqref{eq:plant-ent} 和~\eqref{eq:local-width}，任一条腿均只需要 $O(bu_i+1)$ 比特。对所有块求和，每个辅助插入的这两条腿共需 $O(R)=O(T/a)$ 比特。再对各标记点求和，总共需 $O_{g,\ell}(DT\log D)$ 比特，其中取 $b\ge m$ 以吸收舍入开销。

选择器在网络内部求和，但外层布局必须知道每个标记点处的能量与宽度。将所有块依序分成 $M$ 个等大的连续组。在外层求和中，为每组指定其各数对（类型，$B$ 能量元组）的直方图。直方图记录各数对的重数，私有选择器记录它们的次序。字母表固定，与 $T$ 无关，因此直方图在每个单元内共需 $O_{D,b,M}(\log(T+1))$ 比特。它们决定每组中各条腿的能量和编码连接后的总长度。表将这些直方图与选择器相互核验，并按照选择器给出的块顺序进行解码。每个合法选择器恰好出现于一个外层求和项中。因此，可以先为分组布局选择 $M$，再为选择器布线选择 $\epsilon$，最后为价格比较及选择器宽度选择 $b$，所有选择都在 $T$ 增长之前完成。
